% Propietario conservado: /Users/ruben/Documents/New project/output/REVISION_ARGUMENTAL_APERTURAS_CIERRES_20260915/VI/delivery/MANUSCRIPT_VI_ES_EN_COMPLETE/source_es/sections/14_generacion_regional.tex
\section{Condiciones iniciales y generación regional}
\label{vi:vi:reg:seccion}\label{vi:sec:extension}\label{vi:sec:generacion}

Una posición de APP no contiene por sí sola la condición inicial de su
dinámica. Ésta incluye la dirección de transporte y distingue los cursores
de las dos hojas. Antes de clasificar las partículas es necesario
construir ese dominio: de él proceden los registros regionales que
intervendrán en la coordenada electrónica, la acción y los lectores de
masa. Las operaciones de esta sección no reciben una especie física ni
un valor metrológico como condición inicial.

\subsection{Dos cursores orientados sobre el soporte completo}

Empleamos índices \(I_9=\{0,\ldots,8\}\), cuya etiqueta positiva es
\(i+1\), y las direcciones cardinales \(\mathcal D=\{N,E,S,O\}\).
El dominio de un cursor y el dominio conjunto son
\begin{equation}
 \mathcal S=I_9^2\times\mathcal D,\qquad
 \Omega_{\mathrm{seed}}=\mathcal S_+\times\mathcal S_\times,
 \qquad |\mathcal S|=324,\quad
 |\Omega_{\mathrm{seed}}|=104\,976.
 \label{vi:eq:semillas}
\end{equation}
El subíndice \(\mathrm{seed}\) conserva la denominación del dominio en
los antecedentes; matemáticamente se trata de condiciones iniciales
orientadas. El producto cartesiano incluye todas las parejas ordenadas.
La hoja aditiva dispone de un cursor y la multiplicativa de otro; una
coincidencia posterior entre sus lecturas es un resultado del recorrido.

En los \(54\) instantes de la emisión inicial, el selector de fase
repite \((+1,-1,0)\). En el régimen positivo se evalúa primero
\(\rho_9(i+j+2)\) en el cursor aditivo y se transporta después ese
cursor una posición. En el negativo se evalúa
\(\rho_9((i+1)(j+1))\) en el multiplicativo y después se transporta
el segundo cursor. El régimen nulo registra un evento neutral. Las
direcciones se invierten después de los instantes \(27\) y \(54\).
Esta última inversión pertenece al estado terminal aunque no altera
las lecturas ya efectuadas.

El observable multiplicativo usa la carta de exponentes del grupo de
unidades de \(\ZZ/9\ZZ\):
\[
 \ell_9(1)=0,\quad\ell_9(2)=1,\quad\ell_9(4)=2,\quad
 \ell_9(8)=3,\quad\ell_9(7)=4,\quad\ell_9(5)=5.
\]
En el sector \(\{3,6,9\}\), este observable toma el valor cero.
La pertenencia al sector radical sigue presente en el estado; asignar
ese valor a una lectura escalar no identifica el radical con la unidad
multiplicativa.

\subsection{Fórmula finita de las firmas de ventana}

Las seis ventanas tienen nueve instantes cada una. Cada hoja se evalúa
tres veces por ventana. Una fórmula cerrada permite reconstruir las
firmas sin consultar un listado previo. Para un cursor de origen
\(p=(i,j)\) y dirección \(d\), sea
\[
 f(k)=\begin{cases}k,&0\le k<9,\\18-k,&9\le k<18,
 \end{cases}\qquad p_k=p+f(k)\nu_d\pmod9.
\]
El índice \(k\) cuenta solamente las evaluaciones de esa hoja. La
inversión después de nueve de sus evaluaciones explica la segunda
expresión de \(f\). Escribiendo \(p_k=(i_k,j_k)\), para
\(m=0,\ldots,5\) se obtiene
\begin{align}
 S_m(p,d)&=\sum_{k=3m}^{3m+2}\rho_9(i_k+j_k+2),\notag\\
 E_m(p,d)&=\sum_{k=3m}^{3m+2}
       \ell_9\!\left(\rho_9((i_k+1)(j_k+1))\right),\notag\\
 s_m&=[S_m]_{10},\qquad e_m=[E_m]_{10}.
 \label{vi:vi:reg:firmas}
\end{align}
La evaluación entera y su cociente de división siguen almacenados en
el registro. Las firmas anteriores son las proyecciones decimales
especificadas por el emisor.

Como cada ventana tiene \(Z_m=3\) eventos neutrales, la regla que
acopla ambas firmas es
\begin{align}
 d_{m,1}&=[S_m]_{10},\qquad
 d_{m,2}=[S_m+E_m]_{10},\notag\\
 d_{m,3}&=[3S_m+5E_m+7Z_m]_{10},\notag\\
 U_m&=100s_m+10[s_m+e_m]_{10}+[3s_m+5e_m+1]_{10}.
 \label{vi:eq:emisor-seis}
\end{align}
El término final \(1\) es la reducción de \(7Z_m=21\) módulo
diez. Los coeficientes y el calendario están declarados en la ley
emisora; las magnitudes físicas no aparecen entre sus argumentos.
La proyección ternaria orientada es
\begin{equation}
 \Psi(U_0,\ldots,U_5)=([-U_0]_3,\ldots,[-U_5]_3)
 \in\FF_3^6.
 \label{vi:eq:psi-catalogo}
\end{equation}
Después se aplica la identificación balanceada del TRIT. El signo
menos pertenece a esta carta y debe transportarse al cambiarla.

\begin{proposition}[Factorización y censo de la emisión]
\label{vi:vi:reg:censo}
Las condiciones iniciales producen \(18\) firmas aditivas y \(26\)
multiplicativas. Su acoplamiento es inyectivo y tiene \(468\)
emisiones. La imagen de \(\Psi\) contiene \(243\) palabras en el
ambiente de \(729\) elementos.
\end{proposition}
\begin{proof}
Las centenas de \(U_m\) recuperan \(s_m\). Si \(b_m\) es la
cifra de decenas, \(e_m=[b_m-s_m]_{10}\); la cifra de unidades
comprueba la tercera congruencia. Por ello ninguna pareja distinta
de firmas tiene la misma emisión.

Para obtener el censo se evalúa \eqref{vi:vi:reg:firmas} en
\(i,j=0,\ldots,8\) y las cuatro direcciones. El apéndice
\ref{vi:vi:reg:tablas} contiene todas las firmas resultantes. Cada
firma aditiva tiene \(18\) preimágenes. Entre las multiplicativas,
\(24\) tienen \(8\), una tiene \(24\) y una tiene \(108\).
Estas sumas son respectivamente \(18\cdot18=324\) y
\(24\cdot8+24+108=324\). Puesto que ambos cursores son
independientes, las multiplicidades se multiplican. Resultan
\(432\) emisiones de peso \(144\), \(18\) de peso \(432\) y
\(18\) de peso \(1944\), con suma
\[
 432\cdot144+18\cdot432+18\cdot1944=104\,976.
\]
Se aplican las seis reducciones de \eqref{vi:eq:psi-catalogo} a cada pareja
de las \(18\) y \(26\) firmas publicadas en el apéndice. La
descomposición completa en órbitas allí incluida da
\(38\cdot6+5\cdot3=243\) imágenes diferentes. Todas las
operaciones de esta enumeración son las sumas, productos y residuos
finitos de las fórmulas anteriores.
\end{proof}

La demostración finita no identifica exhaustividad del dominio con
profundidad ilimitada. Las condiciones iniciales se han censado sobre
un soporte fijo. La profundidad se construye por prolongaciones
compatibles que conservan esas condiciones y sus transportes.

\subsection{Selección de regiones mediante fibras orientadas}

Para \(w=(w_1,\ldots,w_6)\), definamos
\[
 r(w)=(w_3,w_1,w_2,w_5,w_6,w_4),\qquad
 s(w)=(w_4,w_5,w_6,w_1,w_2,w_3).
\]
Las identidades \(r^3=s^2=I\), \(srs=r^{-1}\) realizan
\(D_3\). Permutar del mismo modo las seis componentes de \(U\)
conmuta con \(\Psi\). La enumeración anterior comprueba que las
órbitas conservan las multiplicidades de sus fibras.

Sean \(f(w)\) el número de emisiones sobre \(w\), \(m(w)\)
su multiplicidad total y \(\nu(w)=(n_0,n_1,n_2)\) su censo
ternario. La firma aditiva conserva la clase diagonal
\(d=[i+j]_9\) del origen y el sentido \(ES\) o \(NO\).
Escribimos \(\mathcal D(\mathcal O)\) para el conjunto de estas
clases sobre una órbita.

\begin{proposition}[Regiones distinguidas del catálogo]
\label{vi:vi:reg:seleccion}
Existe una única órbita de seis palabras con cinco emisiones por
palabra y multiplicidades \(432+4\cdot144\). Entre las órbitas
de seis palabras que satisfacen
\[
 f(w)=1,\quad m(w)=144\quad(w\in\mathcal O),\qquad
 \mathcal D(\mathcal O)=\{0,3,6\},
\]
el censo \((2,3,1)\) y el censo \((2,2,2)\) seleccionan,
cada uno, una única órbita. La marca \(d=6\) con orientación
\(ES\) determina los representantes
\[
 w^{\rm cl}=010211,\qquad w^{\rm pr}=201101,
 \qquad w^{\rm au}=121200.
\]
\end{proposition}
\begin{proof}
La tabla completa de las \(43\) órbitas en
\ref{vi:vi:reg:tablas} permite aplicar los predicados. La órbita de
clausura es la única con fibra de tamaño cinco y peso \(1008\);
sus cinco pesos son los indicados. La condición diagonal y la
fibra simple de peso \(144\) dejan seis órbitas. Sus censos
seleccionan respectivamente una con \((2,3,1)\) y una con
\((2,2,2)\). Aplicar las seis permutaciones de \(D_3\) y
comprobar la clase diagonal y la orientación en
\eqref{vi:vi:reg:firmas} fija los tres representantes escritos.
La unicidad usa todas las condiciones: si se omite el conjunto de
diagonales, aparecen cuatro órbitas equilibradas de fibra simple
y peso \(144\).
\end{proof}

Las designaciones clausura, propagación y autoescala se aplican
primero a estas regiones. La sección siguiente introduce sus
caracteres de evaluación. En particular, la firma multiplicativa
\(555555\) distingue una emisión transversal de clausura, mientras
que \((5,5)\) es una posición del soporte electrónico. Las dos
construcciones se relacionan mediante operaciones posteriores y no
por identidad de sus cifras.

\subsubsection{Involución especular de las componentes regionales}
\label{vi:mono-ampl:regiones}

Sobre el bloque regional
\(B=(u^{\rm cl},u^{\rm pr},u^{\rm au})\in(\FF_3^6)^3\),
el intercambio es
\[
 \mathfrak cB=(u^{\rm pr},u^{\rm cl},u^{\rm au}).
\]
Su cuadrado es la identidad y mantiene la componente de autoescala.
Si \(q=u^{\rm cl}+u^{\rm pr}+u^{\rm au}\),
\(a=u^{\rm cl}+u^{\rm pr}-u^{\rm au}\) y
\(d=u^{\rm cl}-u^{\rm pr}\), se recupera el bloque mediante
\[
 u^{\rm au}=2(q-a),\quad s=q-u^{\rm au},\quad
 u^{\rm cl}=2(s+d),\quad u^{\rm pr}=2(s-d).
\]
La demostración usa \(2^{-1}=2\) en \(\FF_3\). El intercambio
conserva \(q,a,s\) y cambia el signo de \(d\). Por tanto,
el agregado simétrico y la diferencia orientada conservan juntos
las dos componentes que el agregado aislado no distingue.
Esta involución regional y la inversión de una ruta cardinal
actúan sobre dominios distintos. La componente fija puede
evolucionar durante el transporte aunque sea fija por el intercambio.

\subsection{Calendario, monodromía y memoria}

El calendario finito se compone de doce ventanas nonádicas. Su
sucesión de grupos es
\begin{equation}
 0\longrightarrow C_{12}\xrightarrow{[b]\mapsto[9b]}C_{108}
 \xrightarrow{[t]\mapsto[t]_9}C_9\longrightarrow0.
 \label{vi:eq:extension108}
\end{equation}
Para \(t=9b+r\), \(0\le r<9\), la composición contiene
el cociente de transporte
\[
 (b,r)+(b',r')=
 \left(b+b'+\left\lfloor\frac{r+r'}9\right\rfloor,
       [r+r']_9\right),
\]
reduciendo la primera coordenada módulo doce. La proyección es
exacta porque su núcleo son los múltiplos de nueve. No admite una
sección homomorfa: ésta daría \(C_{108}\cong C_{12}\times C_9\),
pero los exponentes respectivos son \(108\) y \(36\).

La monodromía describe la acción de una vuelta de la base de fase
sobre su fibra. En la dinámica con memoria, la holonomía nonádica
\(\Gamma_9\) conserva la fase y avanza la profundidad. Su
proyección a fase y contador cumple
\[
 (r,n)\longmapsto(r,n+1).
\]
Esta fórmula es una representación reducida de la acción, no una
definición que elimine el resto del estado. Sobre historias
bidireccionales el contador puede extenderse a \(\ZZ\); sobre
historias iniciadas en profundidad cero, el retroceso sólo está
definido donde se conserva el predecesor. El calendario puede
retornar al cabo de \(108\) transiciones y conservar una memoria
distinta de la inicial. Esa distinción será esencial cuando se
construya la partícula como persistencia y no como repetición de
una coordenada visible.

\subsection{De los registros a los operadores de prolongación}

Los operadores lineales de publicación se determinan sobre las
transiciones del transporte. Si \(x_j^\chi\) es un estado regional
y \(\Pi_6\) su publicación de seis componentes, la dependencia es
\[
 b_j^\chi=\Pi_6(x_j^\chi),\qquad
 x_{j+1}^\chi=U_{9j+9}\cdots U_{9j+1}(x_j^\chi).
\]
Las matrices de orígenes \(X_a\) y destinos \(Y_a\), en
vectores fila de \(\FF_3^6\), son
\[
\begin{array}{c|c|c|c}
X_0&Y_0&X_1&Y_1\\\hline
010211&012222&010211&002111\\
012222&010211&002111&110221\\
201101&121221&102011&012222\\
121221&102011&012222&102011\\
121200&112202&121020&010210\\
112202&121020&010210&010200
\end{array}
\]
Se tiene \(\det X_0=1\), \(\det X_1=-1\) en \(\FF_3\).
La solución única de \(X_aL_a=Y_a\) es \(L_a=X_a^{-1}Y_a\).
La concatenación de bloques publicados en el calendario
\((L_0,L_0,L_1,L_1)\) produce longitudes
\(6,12,18,24,30\). La frontera de longitud \(36\) conserva
el terminal y la relación de prolongación. Estas matrices
representan transiciones: la inversión lineal prueba su unicidad
para el registro dado, mientras que la generación de los estados
pertenece a la composición efectiva \(U_t\) del TPK.

\subsection{Alcance de la construcción regional en el estudio de masas}

El censo ha separado condiciones iniciales, emisiones, regiones y
operadores que representan sus transiciones. La siguiente etapa
evalúa las coordenadas regionales y compone el registro dodecafásico.
Más adelante, el atlas de partículas se construirá sobre el mismo
soporte pero empleará otra clasificación. Esta continuidad preserva
una procedencia común sin convertir una fibra numérica en una
especie por simple coincidencia de cardinales o etiquetas.
